\section{Trabalhos Relacionados}

O aprendizado de máquina tem sido aplicado a diferentes decisões do pipeline
variacional. Meta-VQE compartilha informação entre Hamiltonianas parametrizadas
e aprende um conjunto de parâmetros capaz de representar perfis de energia
\cite{cervera2021metavqe}. Essa formulação atua diretamente nos parâmetros do
circuito e explora a continuidade entre instâncias relacionadas.

Outra linha busca pontos iniciais para reduzir o custo da otimização. FLIP
emprega meta-aprendizado para inferir inicializações de circuitos parametrizados
e pode transferir o conhecimento para circuitos maiores que os vistos no treino
\cite{sauvage2021flip}. Qracle representa conjuntamente a Hamiltoniana e o
ansatz como um grafo e usa uma GNN para prever parâmetros iniciais
\cite{zhang2025qracle}. O dataset VQEzy foi proposto para sustentar esse tipo de
aprendizado, reunindo especificações e trajetórias de otimização em diferentes
domínios \cite{zhang2025vqezy}. Esses trabalhos evidenciam tanto o potencial do
aprendizado quanto a importância de dados padronizados.

Há ainda abordagens que alteram a arquitetura do circuito. Um agente de
reinforcement learning pode construir ansätze
incrementalmente, equilibrando acurácia e profundidade no LiH. Essa estratégia
tem espaço de ações e objetivo diferentes do presente estudo: ela sintetiza um
circuito, enquanto aqui se ordena um conjunto discreto e previamente definido
de configurações \cite{ostaszewski2021rl}.

A escolha do otimizador também não pode ser reduzida ao número nominal de
iterações. Um benchmark de 30 otimizadores em 372 instâncias VQE mostrou que
rankings por energia final e por chamadas da função podem divergir, e que
hiperparâmetros internos alteram substancialmente o resultado
\cite{jones2024optimisers}. Isso motiva separar, no presente estudo, redução do
número de configurações candidatas de uma afirmação sobre custo computacional.

Nosso meta-recomendador tampouco prevê a energia FCI, os ângulos do circuito ou
a próxima atualização do otimizador. Ele seleciona, antes da execução, a
combinação de ansatz, otimizador e profundidade com menor erro esperado. Essa
separação permite avaliar a contribuição por métricas de recomendação e pela
fração do grid que deixa de ser executada.
